%
% Front matter
%

\begin{frontmatter}

\title{
  T2K 3.23$\times10^{19}$-POT Muon-Neutrino Disappearance Analysis
}

\vspace*{0.7cm}

\author[RAL]     {Costas Andreopoulos\corref{cor}}
\author[RAL,OXF] {Thomas Dealtry}
\author[ICL,RAL] {James Dobson}
\author[RAL]     {Nick Grant}
\author[OXF]     {Martin Haigh}
\author[RAL]     {Jelena Ilic}
\author[RAL]     {Geoff Pearce}
\author[RAL,OXF] {and Alfons Weber}

\cortext[cor]{E-mail: costas.andreopoulos@stfc.ac.uk}

\address[RAL]{Particle Physics Dept., Rutherford Appleton Laboratory, STFC \\
              Harwell Science and Innovation Campus\\
              Oxfordshire OX11 0QX, UK}
\address[OXF]{Subdepartment of Particle Physics, University of Oxford \\
              Denys Wilkinson Building, Keble Road\\
              Oxfordshire OX1 3RH, UK}
\address[ICL]{Physics Dept., Imperial College London \\
              Blackett Lab, Prince Consort Road\\
              London SW7 2BW, UK}

%
% Abstract
%

\begin{abstract}

We report the results of the Rutherford-Oxford $\nu_{\mu}$-disappearance analysis
on the dataset collected during the first physics run which spanned the period 
from January to June 2010.
This first physics dataset corresponds to an integrated JPARC neutrino beam 
exposure of $3.23 \times 10^{19}$ POT \cite{SuperKAnalysisInputs}.
The oscillation analysis was performed using both GENIE \cite{Genie:NIM} and NEUT \cite{Neut:ActaP}.

Folding in the ND280 measurement, GENIE predicts 22.42 $\pm$ 4.74 (stat) $^{+3.58}_{-3.48}$ (syst)
1-ring $\mu$-like events without oscillations but only 8 were observed.
The statistical significance of the deficit is 2.5 standard deviations.
NEUT predicts 22.81 $\pm$ 4.78 (stat) $^{+3.58}_{-3.49}$ (syst), 
which implies that the statistical significance
of the deficit is 2.5 standard deviations. 

The observed reconstructed energy spectrum of 1-ring $\mu$-like events is consistent
with 2-flavour neutrino oscillations.
A 2-flavour neutrino oscillation fit of the GENIE prediction to the $3.23 \times 10^{19}$ POT dataset, 
with $sin^{2}2\theta_{23}$ constrained in the physical region 
and $|\Delta m^{2}_{23}|$ constrained in a $\sim$15$\sigma$ window around the MINOS best-fit value,
yields 
$|\Delta m^{2}_{23}|$ = 2.8$\times$10$^{-3}$ $eV^{2}/c^{4}$ and
$sin^{2}2\theta_{23}$ = 0.81
as the best-fit values.
For NEUT, the corresponding values are 
$|\Delta m^{2}_{23}|$ = 2.9$\times$10$^{-3}$ $eV^{2}/c^{4}$ and
$sin^{2}2\theta_{23}$ = 0.83.

\end{abstract}

%
% Keywords
%

%\begin{keyword}
%\end{keyword}

\end{frontmatter}

\clearpage

\tableofcontents

\clearpage

\listoffigures
\listoftables


